\section{Realigning and propagating the marked comparison}
\label{sec:allarity-marked-propagation}

The finite polynomial comparison must survive each block integration.
At an intermediate level the marked expansion contains the desired
principal insertion, a small regular remainder, and a forced remainder.
We explain how its leading pieces are reassigned to the prescribed
supports and how the resulting errors accumulate. The identities use
the actual representation at that level.

\subsection{Why a connected principal component stays in one cell}

Let $\rho\ge0$ be an elementary slice radius. Each elementary principal
piece has a nonempty set of marked reader sites; each such reader belongs
to its full support, and the full support is within distance $\rho$ of
every one of those readers. Assume distinct
replica cells have reader separation strictly greater than $2\rho$, and
\[
                  (3+2n)\rho\le dB.
\]
Consider a connected component in the graph joining elementary pieces
whose full supports intersect. Two marked readers on the same piece
are at distance at most $\rho$, since a reader lies in that piece's
support. Readers on two adjacent pieces have distance at most
$2\rho$, by using a point in the intersection. Both bounds are at most
$2\rho$, so the readers belong to the same replica cell. Following a
path of intersecting pieces shows that all marked readers of the
component belong to one cell.

There are at most $n$ distinct marker labels in that cell. Make a graph
on the labels used by the component, joining labels appearing on the
same or on adjacent pieces. The preceding support path gives a walk
between any two of these labels. Erase cycles from the walk. The
remaining simple path uses at most $n-1$ edges, each of geometric
length at most $2\rho$. Thus the distance between any two marked
readers is at most $2n\rho$. Every point in the component's full
support is at distance at most $\rho$ from one of its readers, so
it is at distance at most $(2n+1)\rho\le dB$ from every marked
reader of the component. Its marker diameter is also at most $dB$.
This proves the required enlarged-support bounds independently of
how many elementary pieces the component contains.

For the selected construction the elementary slice radius is $\rho=3B$.
The purchase $6n+3\le a/(3+2n)$ gives the final distance bound
$(2n+1)\rho=(6n+3)B\le aB/(3+2n)$; it does not construct
the slices. In particular it supplies the bound $aB$ required for
the enlarged full supports in the regrouping below. The selected
replica separation is stronger than $2\rho$. Subsequent regrouping of pure principal
components preserves their full supports and marker unions, so the
same argument applies to the principal connected logarithm.

\subsection{Marker prices and spatial prices}

Fix $n\ge1$ marks in each replica cell. Every principal piece has a
nonempty marker subset $I\subseteq\{1,\ldots,n\}$ and a prescribed
spatial support $X_{c,I}$ in cell $c$. The geometric hypotheses are the
following: pieces with the same marker mode have a unique prescribed
support; each actual principal support is contained in that support;
and two prescribed supports containing the same spatial site belong
to the same replica cell. These statements are made for full inspected
supports, including the masks. Let
\[
 \operatorname{chg}_e(X)=\exp(e\operatorname{cov}_B(X)),\qquad
 \operatorname{chg}_e(X_{c,I})\le R_e,
\]
where $e\ge0$ is the support exponent, $B$ the current radius and
$\operatorname{cov}_B$ its covering number. Use one constant marker
weight $\omega>0$ at every mark, of every cell, replica and period.
The weight of a mode $I$ is $\omega^{|I|}$. We do not assume
$\omega\le1$; in the application below $\omega\ge1$. Since
$I\subseteq\{1,\ldots,n\}$, every mode satisfies
\begin{equation}\label{eq:allarity-mode-weight}
 \omega^{|I|}\le\bar\omega_n,\qquad
 \bar\omega_n=\max(1,\omega)^n .
\end{equation}
For $\omega\le1$ this is the bound $\omega^{|I|}\le1$.

If each prescribed piece has weighted coefficient price at most $Z$,
its total incidence price at a fixed mark $(c,j)$ is at most $2^nZ$.
Indeed only subsets $I$ of that same cell can contain $(c,j)$, and
there are at most $2^n$ subsets. This estimate is independent of the
number of replica cells.

A second estimate converts incidence price to spatial price. Suppose
actual pieces $s_a$ have nonnegative prices $p_a$ and nonempty marker
supports, and the total price of pieces containing any fixed mark is
at most $N$. Assign piece $a$ to its prescribed support
$X_{c(a),I(a)}$. At a spatial site $v$, if any such support contains
$v$, all such supports lie in one cell $c$. Charge each summand to
one of its nonempty marks and then overcount by summing over every
mark in that cell. It follows that
\begin{equation}\label{eq:allarity-spatial-incidence}
 \sum_{a:\,v\in X_{c(a),I(a)}}
            p_a\operatorname{chg}_e(X_{c(a),I(a)})
                    \le n N R_e.
\end{equation}
This is why the spatial charge is proportional to $n$, while the
number of principal subset types is bounded by $2^n$.

\subsection{Replacing the leading pieces on their prescribed supports}

Write the current regular source as a sum of remainder pieces $H_a$
and leading pieces $S_b$. Let $F_{c,I}$ be the desired principal source
at mode $(c,I)$, with support $X_{c,I}$. Group the $S_b$ by this marker
mode and define on its healthy region
\[
 L_{c,I}(U)=\sum_{b:\,\operatorname{mode}(S_b)=(c,I)}S_b(U)
                        -F_{c,I}(U).
\]
Assume $|L_{c,I}(U)|\le y$ there. Attach the full support $X_{c,I}$
to $L_{c,I}$ and apply its healthy mask. The incidence price of these
new regular remainders is at most
\begin{equation}\label{eq:allarity-realign-regular}
                         2^n \bar\omega_n R_e y.
\end{equation}
Indeed the new piece at mode $(c,I)$ has value at most $y$ on its
healthy region, support charge at most $R_e$ and marker weight
$\omega^{|I|}$, so its weighted price is at most
$\omega^{|I|}R_ey\le\bar\omega_nR_ey$ by
\eqref{eq:allarity-mode-weight}. The preceding subset count applies.

The mask change also has a forced part. A leading piece may be healthy
on its smaller support while its prescribed larger support is not
healthy. Split its original masked value into its value on the larger
healthy region and the complementary part. The latter is supported
where the enlarged support contains a flag. Insert these complementary
pieces into the forced alphabet, keeping their original marker modes.
The identity is an identity of functions on every field: on the larger
healthy region it is the defining equality for $L_{c,I}$; on its
complement the difference has just been assigned to the forced part.
Apply the exact regular/forced regrouping of
Section~\ref{sec:soft-log-fusion} to the enlarged supports. No division
by an activity value is involved.

Here is its quantitative form. Suppose the previous forced spatial
price is $s$, the leading incidence price is $Z$, and the new forced
inflation exponent is $g\ge0$. If
\begin{equation}\label{eq:allarity-realign-forced-input}
 e^g s+n Z R_{e+1}\le\delta,
           \qquad (2B+1)^4\delta\le1,
\end{equation}
then the regrouped forced price is at most
\begin{equation}\label{eq:allarity-realign-forced-output}
                           e^{-g}(e^\delta-1).
\end{equation}
Here is the proof for the marker-valued factors. Work in the
commutative algebra generated by the marker variables $z_j$ with
$z_j^2=0$. For $I$ a subset of markers write $z_I=\prod_{j\in I}z_j$.
The norm $\|\sum_I a_I z_I\|_\omega=\sum_I |a_I|\omega^{|I|}$ is
submultiplicative for every $\omega>0$: a pair of monomials with
overlapping marker sets multiplies to zero, and otherwise, for
disjoint $I,J$, the weight $\omega^{|I\cup J|}=\omega^{|I|}\omega^{|J|}$
is the product of their weights. No bound $\omega\le1$ is used. Every moved piece is a single $\ell_j=z_{I_j}h_j$ with
$I_j\ne\varnothing$. Thus $\ell_j^2=0$ and $\exp(\ell_j)=1+\ell_j$.
The mask difference defining $h_j$ vanishes whenever the enlarged
support is healthy, so a nonzero moved factor has a flag in its full
support.

Expand the exponential of the sum of the moved pieces as the product
of the factors $1+\ell_j$ and multiply by the old forced gas. In a
nonzero summand each moved factor meets an old forced support: its flag
is covered by the old forced inventories. Taking connected components
of the full-support intersection graph and flattening them back
therefore gives the bijection of
Section~\ref{sec:soft-log-fusion}. This argument preserves the
inventory and the product in the marker algebra. Overlapping marker
sets contribute zero on both sides. The expansion is justified by
absolute convergence in the weighted norm, since the constituent
prices are summable on the finite spatial lattice.

Assign inflated weights $e^g p_i$ to old forced pieces of price $p_i$
and unchanged weights to the moved pieces. At exponent $e+1$ the
site sum is at most $\Gamma=e^g s+nZR_{e+1}$ by
\eqref{eq:allarity-spatial-incidence}. Every surviving connected
component contains an old forced piece, so its product norm is
bounded by $e^{-g}$ times the product of the inflated weights.
The distinct connected-set bound \eqref{eq:ope-forest-price},
with reserve one and counting factor $(2B+1)^4$, bounds the new
site price at exponent $e$ by $e^{-g}(e^\Gamma-1)$.
Since $\Gamma\le\delta$, this proves
\eqref{eq:allarity-realign-forced-output}. The regrouping does not
change the archived background: setting all markers to zero kills
every moved factor, leaving the original background at its original
address.
Thus a regular remainder price $z$ becomes at most
$z+2^n\bar\omega_nR_e y$, by \eqref{eq:allarity-realign-regular},
while the forced price changes as displayed, and the polynomial
density is unchanged.

\subsection{The recurrence along the comparison interval}

Fix a top comparison depth $k$, an intermediate level $l$, a larger
arity allowance $u\ge\max(2,n)$, and a monitor $H>1$ fixed by $k$.
Here $k,l,u,n$ are natural numbers with $l\le k$, and all scalar
price bounds are nonnegative.
Let the tolerance exponent be $q\ge6$ and the regular step multiplier
be $\alpha\ge0$. Set
\[
 d_0=H^{-q},\qquad d_j=(\alpha+2)^j d_0.
\]
The principal sources are the connected polynomials of the same
Gaussian history at the current level, with the same cell and marker
labels. At the terminal comparison level their higher connected
coefficients vanish: the formal logarithm of
$\prod_i(1+z_iP_i)$ is $\sum_i z_iP_i$. Its singleton terms are
exactly the launched calibrated insertions. Consequently the initial
marked representation has regular remainder price $d_0$, after the
specified terminal-source purchase, which must also bound its initial
forced price by $G_k$.

Suppose the actual one-step marked compiler at each level has the
following quantitative outputs when its incoming remainder price is
$z\ge d_0$. Before realignment it supplies remainder price at most
$(\alpha+1)z$, principal incidence price at most $H^{-2}$, and forced
price at most $G_i^2$ at the coarse level. Here
$G_i=e^{-g_i}\le1$ is the required final forced price. Its projected
principal polynomial differs from the next prescribed connected
polynomial by at most $y$ on every required healthy region. Purchase
\begin{gather}
 2^n\bar\omega_nR_e y\le d_0,\qquad
 G_i+n H^{-2}R_{e+1}\le H^{-1},\label{eq:allarity-step-purchases}\\
 (2B_i+1)^4H^{-1}\le1,\qquad e^{H^{-1}}-1\le1.
 \label{eq:allarity-step-small}
\end{gather}
The forced inequality is \eqref{eq:allarity-realign-forced-input}
with $s=G_i^2$, $g=g_i$ and $\delta=H^{-1}$. Its output is at most
$G_i$. The regular price becomes at most
\[
       (\alpha+1)z+d_0\le(\alpha+2)z.
\]
Induction therefore bounds the remainder at level $i$ by $d_{k-i}$,
with forced price $G_i$ and the exact prescribed principal polynomial.
At the induction step the density is the push of the preceding marked
density along the same transition kernel; the realignment changes only
its representation. Hence this is an estimate on the original launched
observable, not a sequence of separately chosen approximants.

The total number of steps is purchased so that
$(\alpha+2)^{k-l}\le H$. It follows that
$d_{k-l}\le H^{1-q}$.

\subsubsection*{The marker weight}
Throughout the comparison interval, at every level from the launch
at the top to the output torus of the replica estimate, the marker
weight is
\begin{equation}\label{eq:allarity-marker-weight}
 \omega=H^{1/(8u)},\qquad H=\mathcal M_k ,
\end{equation}
the same at every mark, replica and finite period. Here
$\mathcal M_k=\exp\{(\log\log\widehat\beta_k)^2\}$ is the clock
monitor at the comparison depth, as in
Section~\ref{sec:rank-root-normalization}. Every price in the
recurrence, namely the remainder prices $d_j$, the forced prices
$G_i$, the principal incidence price $H^{-2}$, the discrepancy
contribution \eqref{eq:allarity-realign-regular} and the terminal
prices, is a price in the $\|\cdot\|_\omega$ norm at this one weight.
The marker weight is never changed inside the interval, so no
rescaling identity between two weighted norms is needed.

Since $H>1$, $\omega\ge1$. As $n\le u$, \eqref{eq:allarity-mode-weight}
gives
\begin{equation}\label{eq:allarity-weight-cap}
 \omega^{|I|}\le\bar\omega_n=\omega^n\le\omega^u=H^{1/8}
 \qquad (I\subseteq\{1,\ldots,n\}).
\end{equation}
Thus the first purchase in \eqref{eq:allarity-step-purchases} is
implied by
\begin{equation}\label{eq:allarity-weighted-discrepancy}
 2^nR_eH^{1/8}y\le d_0=H^{-q}.
\end{equation}
The large weight costs exactly the factor $H^{1/8}$ in the
discrepancy purchase and nothing else in this section: the
incidence and spatial counts, the forced regrouping and the scalar
envelope below hold for every positive weight. The one-step outputs
assumed above, the principal-source bounds and the terminal-source
purchase are hypotheses on prices at the weight
\eqref{eq:allarity-marker-weight}; this section does not derive them
from estimates at a smaller weight. When such an input is available
as an unweighted bound $b$ on the values of pieces carrying at most
$n$ markers, with the support charge already included,
\eqref{eq:allarity-weight-cap} converts it to the weighted price
$H^{1/8}b$. In particular the per-step discrepancy must be supplied
with $y\le2^{-n}R_e^{-1}H^{-q-1/8}$; this is a factor $H^{1/8}$
stronger than the bound $2^nR_ey\le d_0$ that would suffice at a
weight $\omega\le1$.

If the later propagation requires an extra factor $H^{uV+2}$ in the
regular price, purchase $1-q\le-uV-2$. Then
$H^{1-q}\le H^{-uV-2}$ supplies that margin. These two inequalities
are separate purchases: a small initial error does not alone control
arbitrarily many steps or an arbitrary later marker-weight change.

\subsection{The scalar envelope used in the replica estimate}

We record the final elementary step linking a paid marked representation
to Section~\ref{sec:rank-root-normalization}. Suppose a polynomial
density on a torus of $q_0$ sites has the form
\[
 D(U)=e^{\operatorname{exponent}_P(U)}
       \exp_*\left(\sum_a z_{I_a}h_a(U)\right)
       \sum_{B\ {\rm admissible}}\prod_{b\in B}K_b(U).
\]
The $I_a$ are nonempty, and each background packet has nonempty full
spatial support. Let $\|\cdot\|_{\boldsymbol r}$ denote the coefficient
$\ell^1$ norm weighted by positive marker weights. Assume regular
incidence price at most $\ell$ per mark and forced spatial price at
most $L$ per site, with nonnegative support exponent. If there are
$M_0$ marks in total, double counting gives
\[
 \sum_a\|z_{I_a}h_a\|_{\boldsymbol r}\le M_0\ell,
       \qquad\sum_b\|K_b\|_{\boldsymbol r}\le q_0 L.
\]
A nonempty support contributes at least once to the corresponding
incidence sum. Dropping the support weight is legitimate because it
is at least one. Masks can only lower the regular value bounds.
Submultiplicativity gives
\[
 \left\|\exp_*\left(\sum_a z_{I_a}h_a\right)\right\|_{\boldsymbol r}
       \le e^{M_0\ell}.
\]
Dropping the admissibility restrictions in the nonnegative majorant
bounds the forced sum by
$\prod_b(1+\|K_b\|_{\boldsymbol r})\le e^{q_0 L}$. Therefore
\begin{equation}\label{eq:allarity-density-envelope}
 \|D(U)\|_{\boldsymbol r}
      \le\exp\{\operatorname{exponent}_P(U)+q_0L+M_0\ell\}.
\end{equation}
For the replica estimate, first complete any later propagation to the
selected output torus whose site count is exactly the number $q_0$ of
replica cells. This equality is a property of that selected output
level, not of an arbitrary intermediate torus. At this output there
are $M_0=q_0 n$ marks. Thus $L,\ell\le1$
give per-site expense at most $1+n$, exactly the arity dependence
in the determinant ticket. The weight of the coefficient carrying
a chosen nonempty subset $I$ in each of the $q_0$ cells is the
$q_0$th power of its single-cell weight. At the constant weight
\eqref{eq:allarity-marker-weight} this is $\omega^{q_0|I|}$, and
since the coefficient norm dominates each weighted coefficient,
\eqref{eq:allarity-density-envelope} gives pointwise, for the
coefficient $D_{I}^{\otimes}(U)$ of the monomial carrying $I$ in every
cell,
\[
 \omega^{q_0|I|}\,\bigl|D_{I}^{\otimes}(U)\bigr|
   \le\exp\{\operatorname{exponent}_P(U)+q_0(L+n\ell)\}.
\]
With $u=M$ the marker allowance of
Section~\ref{sec:rank-root-normalization}, $\omega=\mathcal
M_k^{1/(8M)}=e^{2I_k}$ is the weight $a$ used there, because
$I_k=\log\mathcal M_k/(16M)$. Hence this is the coefficient price
$a^{q_0|I|}$ with per-site expense $X=L+n\ell\le1+M$ required there,
on the same pushed density. Apply the normalization
estimate and then take the $q_0$th root. The denominator is the
partition function of the same pushed density, preserved by the
unmarked coefficient identity at every step.

The recurrence above states its one-step hypotheses quantitatively.
The joint Gaussian reader estimate supplies the discrepancy $y$ and
the principal source bounds; the mixed compiler supplies the exact
marked push and the pre-realignment prices. The selected scale
construction must buy all of these on the common history, including
the later propagation below $l$. Neither the recurrence nor the
scalar envelope asserts those upstream estimates without their
corresponding purchases.
